\scanpage{62}
\uline{Dém} Soit $\pi$ un cardinal \inserted{infini} tel que les objets de $C$ soient $\pi$-accessibles. Il reste à prouver que la sous-catégorie $E_\pi$ de $E$ formée des objets $\pi$-accessibles est \deleted{\unsure{est}} \inserted{équiv. à une catégorie} \deleted{petite.} \deleted{Il est} Obs $C$ stable par limites inductives finies (sinon on \unsure{s'arrange|s'assure|s’y ramène}, $C$ reste \deleted{\unsure{épuis}} petit) Soit $X$ un objet $\pi$-accessible, on a \deleted{\unsure{lim}} $X=\varinjlim_{C/X}\deleted{\unsure{C/X}}Y$, et $C/X$ est filtrant. \inserted{On a donc aussi $X=\varinjlim_I X_i$, où $I$ est un ens ordonné filtrant, les $X_i\in C$} Soit $J$ l'ens des sous-\deleted{catégories} \inserted{ensembles ordonnés} filtrants \inserted{$I_j$} de $C/X$ de cardinal $\leq\pi$, alors $J$ est \deleted{filtrant et} grand devant $\pi$, et on a
\[
 X=\varinjlim_{j\in J}Y_j\qquad Y_j=\varinjlim_{I_j}\deleted{\unsure{Y_j}}Y_{\unsure{\ell|i}}
\]

Comme $X$ est $\pi$-accessible, \deleted{il faut que l'on ait} l'identité $X\xrightarrow{\sim}\varinjlim_{j\in J}Y_j\simeq X$ se factorise par un $Y_{j'}$, donc $X$ est isom \unsure{à|:} un facteur direct d'un $Y_{j'}$ ! Donc tout objet $\pi$-admissible est isom à un facteur direct d'un objet de la forme $\varinjlim_{i\in I}Y_i$, où \deleted{\unsure{card}} $I$ est un ens \inserted{(\unsure{simplemt|sinon})} ordonné de card. $\leq\pi$, les $Y_i$ dans $C$.
\deleted{\unsure{On trouve} $\Rightarrow$ en particulier \unsure{que} dans chacun filtrant $I\times\mathbf N$ et application \unsure{inductif} donne le résultat \unsure{(+)}} 

\deleted{$Y_{i,n}=Y_i$, $Y_{i,n}\mathbin{-}Y_{j,m}=$ Trouver \deleted{que} ou $\simeq$ au \unsure{lieu} que \deleted{\unsure{si}} pour \uline{tt} cardinal $\pi$ tel que les objets de $C$ soient $\pi$-admissibles, \inserted{$C$ stable par $\varinjlim$ finies,} $E_\pi$ est l'ensemble des facteurs directs des \deleted{systèmes inductifs filtrants} \inserted{limites de} $\varinjlim_I Y_i$, les $Y_i\in C$, $\operatorname{card}I\leq\pi$, \unsure{il suffit de}}

On va en conclure que \deleted{tout \unsure{sys|objet}} $X$ est \deleted{\unsure{ind}} lui-m\textsuperscript{e} une telle limite, ou encore que la catégorie filtrante $C/X$ admet un ensemble cofinal de card $\leq\pi$. \deleted{De façon précise, les $Y_j$ \unsure{$\to X$ cofinaux}}

Je dis que la \deleted{\unsure{classe des}} sous catégorie $F^\pi_{\unsure{\pi|\mathrm R}}$ de $E$ formée des $\varinjlim_{i\in I} Y_{\unsure{i|j}}$, avec $Y_i\in C$, $I$ ens. ordonné tel que $\operatorname{Card} I\leq\pi$, est stable par $\varinjlim$ de cardinal $\leq\pi$. Il suffit de prouver que c'est stable par facteur direct (car on aura prouvé alors que $F^\pi=E_\pi$, et on applique SGA4 I 9.9. Donc il suffit de prouver que c'est stable par $\varinjlim$ filtrantes dénombrables (obvious)\ldots\ldots
% archive sheet number: 61
